\scanpage{106}\mspage{3}
% archive foliation: 105
Si $f\in\mathcal F_3^-$, plus gén. si $f\in\mathcal F_n^-$, alors
$x=-x$, pour $x\in E(k)$, $x\ne0$, implique $f(x)=-f(-x)=-f(x)$
i.e. $2f(x)=0$ i.e. $f(x)=0$. Ainsi $\mathcal F_3\mathbin{-}k^J$ est nul
sur $\mathcal F_3^-$, donc il reste à voir qu'il est
injectif sur $\mathcal F_3^+\simeq\mathcal F_2$. Mais on a mieux

\uline{Corollaire} La matrice \deleted{formée}\inserted{de Vandermonde}
$(\wp(e_i)^\alpha)_{\substack{i\in J\\0\leq\alpha\leq2}}$ est
inversible, en d'autres termes les $\wp(e_i)$ $(i\in J)$
sont tous distincts.
% margin position: left of corollary
\margin{\figure{\begin{tikzpicture}\draw (0,0)--(0,2.5);\end{tikzpicture}}}
\unsure{OPS} \deleted{\unsure{qu'on ait pour pts de}}\inserted{les \unsure{v\textsuperscript{ts}|p\textsuperscript{ts}} de $J$, numérotés $0,1,\infty$,}
et $\Sigma_J=\mathbb P^1_k$. Quitte à multiplier par un
scalaire non nul et ajouter une constante,
\unsure{ops}
\[
\wp(z)=\frac1{z-t}\qquad\text{où }t\unsure{=}k\setminus\{0,1\}=U_{0,3}(k)
\]
donc
\[
\wp(0)=\frac1{-t},\qquad\wp(1)=\frac1{1-t},\qquad\wp(\infty)=0
\]
sont tous distincts, O.K.

\deleted{\unsure{T}} Considérons maintenant
\[
\wp'^2\in\mathcal F_6^+=\deleted{\unsure{\text{munie d'une base}}}
\inserted{\text{espace engendré par les}} \quad1,\wp,\wp^2,\wp^3
\]
et notons que
\[
\varepsilon_6\wp'^2=\varepsilon_6\wp^3=\delta^6
\]
on aura donc la relation
\[
(28)\qquad\boxed{\wp'^2=\wp^3+\unsure{a_1}\wp^2+\unsure{a_2}\wp+\unsure{a_3}}
\]
% margin position: right bottom
\margin{$a_1,a_2,a_3\in k$ déterminés de façon unique
(pour \unsure{$\wp$ choisi})}
